\section{Diferenciabilidad de funciones de \texorpdfstring{$\R$}{R} en \texorpdfstring{$\R$}{R}}
A lo largo de esta sección $f$ denotará una función de valor real definida en un subconjunto  $U\subseteq \mathbb{R}.$
\begin{definicion}
 Sea $c$ un punto interior de $U$. 
 \begin{enumerate}
     \item Diremos que $f$ es diferenciable en $c$ si existe el siguiente límite, el cual llamaremos la derivada de $f$ en $c$ y denotaremos
 \begin{equation}\label{def-derivada}
 f'(c):=\lim_{x\to c}\frac{f(x)-f(c)}{x-a}.    
 \end{equation}
 
 \item Si $U$ es un abierto y $f$ es diferenciable en $a$ para todo $c\in U$, decimos que $f$ es diferenciable (en $U$). La función $f':U\to \mathbb{R}$, se llama la  derivada de $f$.
 \item Otras notaciones para la derivada de $f$  son $\frac{df}{dx}$, $Df$ y $f^{(1)}$.
 \end{enumerate}
\end{definicion}

\begin{remark}
De la ecuación \eqref{def-derivada} tenemos que 
\begin{eqnarray*}
 & \lim_{x\to c}\frac{f(x)-f(c)}{x-c}-f'(c)=0 \\
 &
 \lim_{x\to c}\frac{f(x)-f(c)-f'(c)(x-c)}{x-c}=0.   
 \end{eqnarray*}
\end{remark}

\begin{definicion}
Sea $c$ un punto interior de $U$ y supongamos que $f$ es diferenciable en $c$. La recta tangente a la gráfica de $f$ en el punto $(c, f(c))$ es la recta  cuya ecuación es 
 \begin{equation}\label{def-tangente}
 y=f(c)+f'(c)(x-c) .
 \end{equation}
\end{definicion}
\begin{teorema}
    Si $f$ está definida en $(a, b)$ y es diferenciable en $c \in (a, b)$, entonces existe una función $f^*$ (que depende de $f$ y de $c$) continua en $c$ que satisface:
$$f(x) - f(c) = (x - c) \cdot f^*(x) \quad \forall x \in (a, b), \text{ con } f^*(c) = f'(c)$$
\textbf{Recíprocamente:} si existe tal $f^*$ continua en $c$ que satisfaga la ecuación, entonces $f$ es diferenciable en $c$ y $f'(c) = f^*(c)$.

\end{teorema}
\begin{prueba}
    Sea $x \neq c$ y $f$ es diferenciable.Definimos $f^{*}$ como:
$$f^*(x) = \frac{f(x) - f(c)}{x - c} \quad \text{si } x \neq c, \ f^*(c) = f'(c)$$
Multiplicamos por $(x - c)$ la ecuación anterior,
$$(x - c) \cdot f^*(x) = f(x) - f(c)$$
Si $x = c$, claramente es igual a cero, pues $f(c) - f(c) = 0 = (c - c) \cdot f^*(c)$. Por lo tanto se tiene que $f^*(x)(x - c) = f(x) - f(c)$ para todo $x \in (a, b)$.
Ahora debemos demostrar que $f^*$ es continua en $c$.

Para $x \neq c$:
$$\lim_{x \to c} f^*(x) = \lim_{x \to c} \frac{f(x) - f(c)}{x - c}$$
Pero por hipótesis se tiene que $f$ es diferenciable en $c$, por lo que
$$\lim_{x \to c} f^*(x) = f^*(c).$$
Se concluye que $f^*$ es continua en $c$.

\end{prueba}
\begin{teorema}
Si $f$ es diferenciable en $c$ entonces $f$ es continua en $c$.
\end{teorema}
\begin{prueba}
    Por anterior, tenemos que $f^*$ es continua en $c$, y satisface:
$$f^*(x)(x-c) = f(x)-f(c) \quad \forall x \neq c$$
Ya que $f^*$ es continua en $c$ y $f^*(c) = f'(c)$.
Ahora si $x \to c$ entonces
$$\lim_{x \to c} (f(x)-f(c)) = \lim_{x \to c} (f^*(x)(x-c))$$
Pero como $x \to c$ se tiene que $(c-c) = 0$ y $\lim_{x \to c} f^*(x) = f'(c)$, por lo que:
$$\lim_{x \to c} (f(x)-f(c)) = f'(c) \cdot 0$$
$$\lim_{x \to c} (f(x) - f(c)) = 0$$
Ya que $\lim_{x \to c} f(c) = f(c)$
$$\lim_{x \to c} f(x) = f(c)$$
Por lo tanto se tiene que $f$ es continua en $c$.

\end{prueba}

\begin{example}
    {¿Si $f$ es continua entonces $f$ es diferenciable?}\\
Sea $f(x) = |x|$, analizaremos que es continua y no diferenciable en $0$:
1) $f$ es continua:
$$\lim_{x \to 0} f(x) = f(0)$$
Analizaremos si $\lim_{x \to 0^+} f(x) = \lim_{x \to 0^-} f(x) = f(0)$
\begin{itemize}
    \item $x > 0$ por lo que $|x| = x$ así que:
    $$\lim_{x \to 0^+} |x| = \lim_{x \to 0^+} x = 0$$
    
    \item Si $x < 0$ por lo que $|x| = -x$ así que:
    $$\lim_{x \to 0^-} |x| = \lim_{x \to 0^-} -x = 0$$
\end{itemize}

2) ¿Ahora si $f$ es diferenciable en cero?
Ahora apliquemos la definición de diferenciabilidad:
$$\lim_{x \to 0} \frac{f(x) - f(0)}{x - 0}$$
$$\lim_{x \to 0} \frac{f(x)}{x} = \lim_{x \to 0} \frac{|x|}{x}$$
Ahora analizaremos por los laterales:
\begin{itemize}
    \item $x > 0$, entonces $|x| = x$ por lo que se tiene que:
    $$\lim_{x \to 0^+} \frac{|x|}{x} = \lim_{x \to 0^+} \frac{x}{x} = 1 \quad \text{(1)}$$
    
    \item Si $x < 0$, entonces $|x| = -x$, así que:
    $$\lim_{x \to 0^-} \frac{|x|}{x} = \lim_{x \to 0^-} \frac{-x}{x} = -1 \quad \text{(2)}$$
\end{itemize}
Concluimos que $f(x)$ es continua pero no diferenciable ya que $\text{(1)} \neq \text{(2)}$.

\end{example}
\begin{example}
 Si $f$ es diferenciable en un abierto $U$, no necesariamente $f'$ es continua en $U$. Ilustremos esta afirmación con un ejemplo. Sea $f:\mathbb{R} \to \mathbb{R}$ definida por 
 \[f(x)= \left\{ \begin{array}{rll}
    x^2\sin\left(\frac{1}{x}\right), & \text{si}  & x\neq 0  \\
    0, & \text{si} & x=0 
 \end{array}\right. .\]
Enonces $f$ es diferenciable y 
\[f'(x)= \left\{ \begin{array}{rll}
    2x\sin\left(\frac{1}{x}\right)-\cos\left(\frac{1}{x}\right), & \text{si}  & x\neq 0  \\
    0, & \text{si} & x=0 
 \end{array}\right. .\]
Pero notemos que $f'$ no es continua en $0$. Para ver esto, considere la sucesión $x_n=\frac{1}{2\pi n}$. Entonces es claro que $x_n \to 0$. Sin embargo $f'(x_n)=-1 \to -1 \neq f'(0)=0$.
\end{example}

\begin{definicion}

\begin{enumerate}
    \item Supongamos que $f$ es diferenciable y que $f'$ es continua. Entonces decimos que $f$ es una función de clase 1. Al conjunto de todas las funciones de clase 1, definidas en $U$ lo denotamos $C^1(U)$.
    \item Si $f$ tiene derivada de orden $n$, $n=1,2,\dots $ y $f^{(n)}$ es continua, decimos que $f$ es una función de clase $n$. Al conjunto de tales funciones lo denotamos $C^{n}(U)$.
    \item Denotemos $C^{0}(U)$ al conjunto de todas las funciones continuas definidas en $U$.
\end{enumerate}

\begin{remark}
Notemos que $C^0(U)\supseteq C^1(U)\supseteq C^2(U)\supseteq \cdots .$
\end{remark}
\end{definicion}

\begin{teorema}{Álgebra de Derivadas}
Supongamos que $f$ y $g$ están definidas en $(a, b)$ y son diferenciables en $c$. Entonces $f + g$, $f - g$, $f \cdot g$ son también diferenciables en $c$. Lo mismo es verdadero para $\frac{f}{g}$ si $g(c) \neq 0$.

\begin{enumerate}[label=\alph*)]
    \item \textbf{Suma / Diferencia}
    $$(f \pm g)'(c) = f'(c) \pm g'(c)$$
    
    \item \textbf{Producto}
    $$(f \cdot g)'(c) = f(c)g'(c) + f'(c)g(c)$$
    
    \item \textbf{Cociente $[g(c) \neq 0]$}
    $$\left(\frac{f}{g}\right)'(c) = \frac{[g(c)f'(c) - f(c)g'(c)]}{g(c)^2}$$
\end{enumerate}
\end{teorema}
\begin{prueba}
    \begin{itemize}
        \item  $(f \pm g)'(x) = f'(x) \pm g'(x)$:
Tenemos que $f$ y $g$ son diferenciables en $c$ aplicando el teorema anterior tenemos que: $f(x) = f(c) + (x-c)f^*(x)$ y $g(x) = g(c) + (x-c)g^*(x)$, entonces:
$$(f \pm g)(x) = f(x) \pm g(x)$$
$$(f \pm g)(x) = (f(c) + (x-c)f^*(x)) \pm (g(c) + (x-c)g^*(x))$$
$$(f \pm g)(x) = f(c) \pm g(c) + (x-c)(f^*(x) \pm g^*(x))$$
Ahora aplicamos la definición de la derivada
$$(f \pm g)'(c) = \lim_{x \to c} \frac{[(f \pm g)(c) + (x-c)(f^*(x) \pm g^*(x))] - (f \pm g)(c)}{x-c}$$
$$(f \pm g)'(c) = \lim_{x \to c} (f^*(x) \pm g^*(x))$$
$$(f \pm g)'(c) = \lim_{x \to c} f^*(c) \pm \lim_{x \to c} g^*(c) \quad \text{Propiedad del límite}$$
Como $\lim_{x \to c} f^*(x) = f'(c)$ y $\lim_{x \to c} g^*(x) = g'(c)$, por lo que obtenemos:
$$(f \pm g)'(c) = f'(c) \pm g'(c)$$
\item $(f \cdot g)'(x) = f'(x) \cdot g(x) + f(x) \cdot g'(x)$
Tenemos que $f$ y $g$ son diferenciables en $c$, aplicando el teorema anterior tenemos que $f(x) = f(c) + (x-c)f^*(x)$ y $g(x) = g(c) + (x-c)g^*(x)$, entonces:
$$(f \cdot g)(x) = f(x) \cdot g(x)$$
$$(f \cdot g)(x) = (f(c) + (x-c)f^*(x)) \cdot (g(c) + (x-c)g^*(x))$$
$$(f \cdot g)(x) = f(c) \cdot g(c) + (x-c)g^*(x) \cdot f(c) + (x-c)f^*(x) \cdot g(c) + (x-c)^2 g^*(x) \cdot f^*(x)$$
$$(f \cdot g)(x) = f(c) \cdot g(c) + (x-c)(g^*(x) \cdot f(c) + f^*(x) \cdot g(c) + (x-c)g^*(x) \cdot f^*(x))$$
Aplicando la definición de la derivada:
$$(f \cdot g)'(c) = \lim_{x \to c} \frac{[(f \cdot g)(c) + (x-c)(g^*(x) \cdot f(c) + f^*(x) \cdot g(c) + (x-c)g^*(x) \cdot f^*(x))] - (f \cdot g)(c)}{x-c}$$
$$(f \cdot g)'(c) = \lim_{x \to c} \frac{(x-c)(g^*(x) \cdot f(c) + f^*(x) \cdot g(c) + (x-c)g^*(x) \cdot f^*(x))}{x-c}$$
$$(f \cdot g)'(c) = \lim_{x \to c} g^*(x) \cdot f(c) + f^*(x) \cdot g(c) + (x-c)g^*(x) \cdot f^*(x)$$
$$(f \cdot g)'(c) = g(c) \lim_{x \to c} f^*(x) + f(c) \lim_{x \to c} g^*(x) + \lim_{x \to c} (x-c)g^*(x) \cdot f^*(x)$$
Como $x \to c$ se tiene que $x-c = 0$ por lo que el último término se hace cero y $\lim_{x \to c} f^*(x) = f'(c)$ y $\lim_{x \to c} g^*(x) = g'(c)$ se tiene que:
$$(f \cdot g)'(c) = f'(c) \cdot g(c) + g'(c) \cdot f(c)$$

\item $\left(\frac{f}{g}\right)(x) = \frac{f'(x)g(x) - g'(x)f(x)}{(g(x))^2}$
Como $f$ y $g$ son diferenciables en $c$ y $g(c) \neq 0$ tal que $f(x) = f(c) + (x-c)f^*(x)$ y $g(x) = g(c) + (x-c)g^*(x)$, entonces:
$$\left(\frac{f}{g}\right)(x) = \frac{f(c) + (x-c)f^*(x)}{g(c) + (x-c)g^*(x)}$$
Aplicando la definición de la derivada:
$$\left(\frac{f}{g}\right)'(c) = \lim_{x \to c} \frac{\frac{f(c) + (x-c)f^*(x)}{g(c) + (x-c)g^*(x)} - \frac{f(c)}{g(c)}}{x-c}$$
$$\left(\frac{f}{g}\right)'(c) = \lim_{x \to c} \frac{\frac{f(c) \cdot g(c) + (x-c)f^*(x) \cdot g(c) - f(c) \cdot g(c) - (x-c)g^*(x) \cdot f(c)}{g^2(c) + (x-c)g^*(x) \cdot g(c)}}{x-c}$$
$$\left(\frac{f}{g}\right)'(c) = \lim_{x \to c} \frac{\frac{(x-c)(f^*(x) \cdot g(c) - g^*(x) \cdot f(c))}{g^2(c) + (x-c)g^*(x) \cdot g(c)}}{x-c}$$
$$\left(\frac{f}{g}\right)'(c) = \lim_{x \to c} \frac{f^*(x) \cdot g(c) - g^*(x) \cdot f(c)}{g^2(c) + (x-c)g^*(x) \cdot g(c)}$$
Como $x \to c$ se tiene que $(x-c)g^*(x) \cdot g(c) = 0$ por lo que el último término se hace cero y $\lim_{x \to c} f^*(x) = f'(c)$ y $\lim_{x \to c} g^*(x) = g'(c)$ se tiene que:
$$\left(\frac{f}{g}\right)'(c) = \frac{f'(c) \cdot g(c) - g'(c) \cdot f(c)}{g^2(c)}$$


    \end{itemize}
\end{prueba}
\begin{teorema}{La Regla de la Cadena}
Sea $f$ definida en un intervalo abierto $S$, y sea $g$ definida en $f(S)$. Consideremos la función compuesta $g \circ f$ definida en $S$ por:
$$(g \circ f)(x) = g(f(x))$$
Si $f$ es diferenciable en $c \in S$ y $g$ es diferenciable en $f(c)$, entonces $g \circ f$ es diferenciable en $c$ y
$$(g \circ f)'(c) = g'(f(c)) \cdot f'(c)$$
\end{teorema}
\begin{prueba}
    Tenemos que $f$ y $g$ son diferenciables, aplicamos el teorema 1:
$$f^*(x)(x-c) = f(x)-f(c) \quad \forall x \in S \quad \text{(1)}$$
Sabemos que $f^*$ es continua en $c$ ya que $f^*(c) = f'(c)$.
Ahora sea $d = f(c)$. Como $g$ es diferenciable en $d$, aplicando el teorema 1 en $g$:
$$g(y) - g(d) = (y-d) g^*(y) \quad \forall y \in f(S) \quad \text{(2)}$$
Donde $g^*$ es continua en $d$ y $g^*(d) = g'(d)$.
Sea $y \in f(S)$, existe $x \in S$ tal que $y = f(x)$ entonces de (2) se tiene que:
$$(g \circ f)(x) - (g \circ f)(c) = g(f(x)) - g(f(c)) = (f(x) - f(c)) g^*(f(x))$$
Reemplazamos $f(x) - f(c) = (x-c) f^*(x)$ por (1), por lo que:
$$(g \circ f)(x) - (g \circ f)(c) = [(x-c) f^*(x)] g^*(f(x))$$
$$(g \circ f)(x) - (g \circ f)(c) = (x-c) [f^*(x) g^*(f(x))]$$
Ahora definamos que $h(x) = f^*(x) g^*(f(x)) \quad \forall x \in S$, $h$ es continua en $c$:

\begin{itemize}
    \item $f^*$ es diferenciable en $c$ por el teorema 1
    
    \item $f$ es diferenciable en $c$, entonces $f$ es continua en $c$ por el teorema 2
    
    \item $g^*$ es continua en $d = f(c)$ por el teorema 1
\end{itemize}

Como $g^*$ es continua en $d$ y $f$ es continua en $c$, por la composición $g^*(f(x))$ es continua en $c$, entonces se tiene que:
$$h(x) = f^*(x) g^*(f(x)) \text{ es continua en } c$$
Dado que
$$\lim_{x \to c} f^*(x) = f^*(c) \quad \text{y} \quad \lim_{y \to d} g^*(y) = g^*(d)$$
Entonces $h(c) = f^*(c)g^*(d) = f'(c)g'(f(c))$, ya que $f^*(c) = f'(c)$ y $g^*(d) = g^*(f(c)) = g'(f(c))$.
Como $h$ es continua en $c$, se tiene:
$$\frac{(g \circ f)(x) - (g \circ f)(c)}{x - c} = h(x)$$
$$\lim_{x \to c} \frac{(g \circ f)(x) - (g \circ f)(c)}{x - c} = \lim_{x \to c} h(x) = h(c) = g^*(f(c)) f^*(c) = g'(f(c)) f'(c)$$
Por el teorema 1 se concluye que:
$$(g \circ f)'(c) = h(c) = g'(f(c)) f'(c)$$
Por lo que $(g \circ f)$ es diferenciable en $c$.

\end{prueba}

\begin{definicion}
   Sea $f$ una función real definida en un subconjunto $S$ de un espacio métrico $M$, y $a \in S$.

\begin{itemize}
    \item \textbf{$f$ posee un máximo local en $a$} si existe una bola $B(a)$ tal que:
    $$f(x) \le f(a) \quad \forall x \in B(a) \cap S$$
    
    \item \textbf{$f$ posee un mínimo local en $a$} si existe una bola $B(a)$ tal que:
    $$f(x) \ge f(a) \quad \forall x \in B(a) \cap S$$
\end{itemize}
\end{definicion}
\begin{teorema}
    Si $f$ posee un máximo o mínimo local en un punto interior $c$ de $(a, b)$ y posee derivada finita en $c$, entonces:
$$f'(c) = 0$$
\textbf{Nota sobre derivadas infinitas:} Si la derivada es infinita en $c$ (es decir, existe una tangente vertical), $c$ \textbf{sí puede} ser un extremo relativo.
\end{teorema}
\begin{prueba}
    \textbf{Caso 1:} Supongamos que $f$ posee derivada en $c$ y tiene un máximo local en $c$.
Por definición, existe $\delta > 0$ tal que para todo $x \in (c-\delta, c+\delta)$:
$$f(x) \le f(a)$$
Es decir:
$$f(x) - f(c) \le 0$$

\begin{itemize}
    \item \textbf{Si $x > c$ (por derecha):}
    $$f'(c) = \lim_{x \to c^+} \frac{f(x) - f(c)}{x-c} \le 0 \quad \text{ya que } f(x)-f(c) \le 0 \text{ y } (x-c) > 0$$
    
    \item \textbf{Si $x < c$ (por izquierda):}
    $$f'(c) = \lim_{x \to c^-} \frac{f(x) - f(c)}{x-c} \ge 0 \quad \text{ya que } f(x)-f(c) \le 0 \text{ y } (x-c) < 0$$
\end{itemize}

Tenemos que $f'(c) \le 0$ y $f'(c) \ge 0$, simultáneamente.
Por lo que la única posibilidad es:
$$f'(c) = 0$$
Si $f'(c) = +\infty$, entonces por definición se tiene que:
$$\lim_{x \to c} \frac{f(x) - f(c)}{x - c} = +\infty \iff \forall M > 0, \exists \delta > 0 \text{ tal que si } x \in B(c, \delta) \cap S \implies \frac{f(x) - f(c)}{x - c} > M$$
En particular, con $M = 1$, $\exists \delta > 0$ tal que si $x \in B(c, \delta) \cap S \implies \frac{f(x) - f(c)}{x - c} > 1$.
Si $x > c$:
$$f(x) - f(c) > x - c > 0 \implies f(x) > f(c) \quad \text{(2)}$$
Por lo anterior se tiene que $c$ no es un extremo relativo.

\textbf{Caso 2:} Es análogo al caso 1, pero suponiendo que $f$ tiene un mínimo local en $c$.
En ambos casos (máximo o mínimo local en un punto interior donde existe derivada), obtenemos que:
$$f'(c) = 0$$
\end{prueba}
